% !TEX root = ../thesis.tex
\mychapter{Literature}{Literature}{}
\label{chap:literature}

\section{Related Work/Simulation Based approaches}
\label{sec:relatedwork}


Finding Optimal Camera Placement otherwise known as the photogrammetry network design problem has been a problem people have tried to solve in multiple-contexts.
They have ranged from optimal camera placement for industrial pose estimation to 
There are numerous works done that utilise simulation based approaches

jorgemsem did work on simulation based pose estimation for industrial. This used a lot of algorithms but wasz for tray optimssed not humans 
zhao a 2013 paper sought 
rahami,
olague

then 2025 yang 

Finding optimal camera placement has been a long studied problem for various contexts. Historically the problem of camera placement known as photogrammetry network design 
They have found use cases for cases in industrial pose estimation, reconstruction of points 

Simulation based approaches to solving the Optimal Camera placement problem, known as photogrammetry network design, has been employed for decades \cite{fraser1984} \cite{masson1995} and \cite{olague2001applications}. They were originally classified into four orders of design zero, first, second and third respectively classified as the datum, configuration, weight and densification problems as originally classified by \cite{Grafarend1974} and has been implemented by 
It has been noted that for an optimal multi-camera setup there ought to be focus on accuracy specifications, essentially a measure of correctness, quality of the configuration defined by precision and realiability. precision is a function of the geometry of the camera network.
By discussing and using the classification system, \citet{fraser} describes a general process by which a camera network can be optimised through computer simulation.
Althopu
\citet{masson1995} states that for a high accuracy industrial photogrammetry measurements that an appropriate multi-camera network is required. 

Experts have realied upon simulation based approaches to oovercome the complexity of the camera network design \cite{mason1995}
Aspects that influence the accuracy of camera configuration considered are triangulability, where a point of interest is viewed by more than one camera, the number and position of cameras in the network and lastly the convergence and intersection of cameras for a point of interest. These all contribute to the accuracy of a multi camera network.
\cite{TarboxandGottschlich1993} proved that the FOD problem is NP-complete meaning that it's algorithmic complexity exponential (greater than polynomial). This means that trial and error approaches, ie. bruteforcing the solution space is infeasible and a heuristic approach is required.
\citet{Mason} employs the general formulation by \cite{Fraser1984} 
Following the choice of camera (observation selection scheme), lens, sensor the photogrammetry pracitioner/designer constructs an approximation of a suitable geometry where the bundle method is used. The simulation strategies tackle the orders of design as by \cite{Grafarend1974} but does not answer questions of number of cameras.

\cite{Olague} is revolutionary as it poses the problem in terms of a global optimisation design exploring the solution space by means of both combinatorial search and non-continuous optimisation using a Genetic Algorithm and define the criteria 



\citet{zhao2013approximate,rahimian2016optimal, jorgensen2018simulation}
\section{Motion Capture for biomechanics}

\section{Human Pose Estimation}
\subsection{Top Down Bottom Uo}
\subsection{Multi-view 3D reconstruction}
\subsection{Evaluation Metrics}

\section{Camera placement and camera network design}
\cite{kwak2025assessing}

\cite{rahimian2016optimal}
\cite{zhao2013approximate}
\cite{jorgensen2018simulation}
\cite{yang2025comparison}
\cite{nishikawa2025comparison}
\cite{kim2026comparison}
\cite{wren2023comparison}

\subsection{Human Pose Estimation Algorithms}


\section{Synthetic Data for Human Pose Estimation}

\subsection{Synthetic Dataset}

\cite{varol2017learning}
\cite{patel2021agora}
\cite{yang2023synbody}
\cite{black2023bedlam}
\cite{tesch2025bedlam2}

\subsubsection{Modeling the human body historical}

\subsubsection{Modern advancements}
\cite{loper2015smpl}
\cite{SMPL-X:2019}
\cite{bregier2026human}

\subsection{Motion Data}
\cite{mahmood2019amass}
\cite{ghorbani2021soma}

\subsection{digital content creation tools}
Maybe a table showing the languages, add-ons, rendering, api's

\subsubsection{Unreal}
\subsubsection{Unity}
\subsubsection{Blender}

\subsection{Sim2Real}


\section{Optimisation for combinatorial subset selection problem}

All optimisation methods
 \cite{jorgensen2018simulation} \cite{olague2001applications} \cite{zhao2013approximate}, \cite{rahimian2016optimal}

 \cite{eberhart1995particle}
 \cite{storn1997differential}

\cite{van2023set}
\cite{langeveld2015set}
\cite{khan2012fuzzy}
\cite{mohiuddin2016fuzzy}

\section{Chapter Summary}
